\documentclass[12pt, a4paper]{ctexart}
\usepackage[margin=2cm]{geometry}
\usepackage{fancyhdr}
\usepackage{amsmath,amssymb,amsthm}
\usepackage{enumitem}
\usepackage{graphicx}
\usepackage{tikz}
\usepackage{pgfplots}
\pgfplotsset{compat=1.18}

\setCJKmainfont{Hiragino Sans GB}
\setCJKsansfont{Heiti SC}

\pagestyle{fancy}
\fancyhf{}
\fancyhead[L]{高考数学满分突破·导数深层理解}
\fancyhead[R]{测试卷}
\fancyfoot[C]{\thepage}
\renewcommand{\headrulewidth}{0.4pt}

\newcommand{\blank}[1]{\underline{\hspace{#1}}}

\begin{document}

\begin{center}
{\LARGE\bfseries 高考数学满分突破·导数深层理解\quad 测试卷}\\[0.5em]
{\large 满分150分\quad 考试时间120分钟}
\end{center}

\vspace{1em}

\noindent\textbf{一、单项选择题（本大题共8小题，每小题5分，共40分。每小题只有一个选项符合题目要求。）}

\begin{enumerate}
\item 函数$f(x)=x^3-3x^2+4$在区间$[-1,3]$上的最大值为\blank{3cm}。
\begin{enumerate}[label=\Alph*.]
\item $0$
\item $4$
\item $8$
\item $12$
\end{enumerate}

\item 设$f(x)=xe^{-x}$，则$f(x)$的单调递增区间为\blank{3cm}。
\begin{enumerate}[label=\Alph*.]
\item $(-\infty,1)$
\item $(1,+\infty)$
\item $(-\infty,0)$
\item $(0,+\infty)$
\end{enumerate}

\item 曲线$y=x^3-3x^2+2x+1$在点$(1,1)$处的切线方程为\blank{3cm}。
\begin{enumerate}[label=\Alph*.]
\item $y=-x+2$
\item $y=x$
\item $y=-2x+3$
\item $y=2x-1$
\end{enumerate}

\item 函数$f(x)=x-\ln x$的极小值为\blank{3cm}。
\begin{enumerate}[label=\Alph*.]
\item $0$
\item $1$
\item $-1$
\item $e-1$
\end{enumerate}

\item 已知$f(x)=ax^3+bx^2+cx+d$在$x=-1$处取极大值$4$，在$x=1$处取极小值$0$，则$a+b+c+d$的值为\blank{3cm}。
\begin{enumerate}[label=\Alph*.]
\item $1$
\item $2$
\item $3$
\item $4$
\end{enumerate}

\item 设$f(x)=\dfrac{\ln x}{x}$，则$f(x)$在$[1,e]$上的最大值为\blank{3cm}。
\begin{enumerate}[label=\Alph*.]
\item $0$
\item $\dfrac{1}{e}$
\item $1$
\item $\dfrac{2}{e}$
\end{enumerate}

\item 若函数$f(x)=x^3-12x+a$有三个不同的零点，则$a$的取值范围是\blank{3cm}。
\begin{enumerate}[label=\Alph*.]
\item $(-16,16)$
\item $(-8,8)$
\item $(-4,4)$
\item $(-\infty,-16)\cup(16,+\infty)$
\end{enumerate}

\item 已知$f(x)=e^x-x-1$，则当$x\geq 0$时，$f(x)$与$g(x)=\dfrac{x^2}{2}$的大小关系为\blank{3cm}。
\begin{enumerate}[label=\Alph*.]
\item $f(x)\geq g(x)$
\item $f(x)\leq g(x)$
\item $f(x)=g(x)$
\item 无法确定
\end{enumerate}
\end{enumerate}

\vspace{1em}

\noindent\textbf{二、多项选择题（本大题共4小题，每小题5分，共20分。每小题有多个选项符合题目要求，全部选对得5分，部分选对得2分，有选错的得0分。）}

\begin{enumerate}[start=9]
\item 关于函数$f(x)=x^3-3x$，下列结论正确的有\blank{3cm}。
\begin{enumerate}[label=\Alph*.]
\item $f(x)$在$(-1,1)$上单调递减
\item $f(x)$在$x=-1$处取得极大值$2$
\item $f(x)$在$x=1$处取得极小值$-2$
\item $f(x)$有三个零点
\end{enumerate}

\item 下列不等式恒成立的有\blank{3cm}。
\begin{enumerate}[label=\Alph*.]
\item $e^x\geq x+1$（$x\in\mathbb{R}$）
\item $\ln x\leq x-1$（$x>0$）
\item $\sin x < x$（$x>0$）
\item $x-\dfrac{x^3}{6}\leq\sin x\leq x$（$x\geq 0$）
\end{enumerate}

\item 设$f(x)=\dfrac{x}{e^x}$，则下列结论正确的有\blank{3cm}。
\begin{enumerate}[label=\Alph*.]
\item $f(x)$在$x=1$处取得最大值$\dfrac{1}{e}$
\item $f(x)$在$(-\infty,1)$上单调递增
\item 当$x>e$时，$f(x)<\dfrac{1}{e}$
\item 方程$f(x)=a$（$a>0$）最多有两个解
\end{enumerate}

\item 关于导数与原函数的关系，下列说法正确的有\blank{3cm}。
\begin{enumerate}[label=\Alph*.]
\item 若$f'(x_0)=0$，则$x_0$是$f(x)$的极值点
\item 若$f'(x)>0$在$(a,b)$上恒成立，则$f(x)$在$(a,b)$上单调递增
\item 若$f''(x_0)=0$且$f'''(x_0)\neq 0$，则$(x_0,f(x_0))$是拐点
\item 若$f(x)$在$[a,b]$上连续，在$(a,b)$上可导，且$f(a)=f(b)$，则存在$\xi\in(a,b)$使得$f'(\xi)=0$
\end{enumerate}
\end{enumerate}

\vspace{1em}

\noindent\textbf{三、填空题（本大题共4小题，每小题5分，共20分。）}

\begin{enumerate}[start=13]
\item 函数$f(x)=x^3-3x+1$在$[-2,2]$上的最小值为\blank{3cm}。

\item 已知$f(x)=\ln x-ax$（$a>0$），若$f(x)\leq -1$恒成立，则$a$的值为\blank{3cm}。

\item 曲线$y=e^x$与$y=\ln x$的公切线方程为\blank{4cm}。

\item 设$f(x)=x^2\ln x$，则$f(x)$的单调递减区间为\blank{3cm}。
\end{enumerate}

\vspace{1em}

\noindent\textbf{四、解答题（本大题共6小题，共70分。解答应写出文字说明、证明过程或演算步骤。）}

\begin{enumerate}[start=17]
\item （10分）已知函数$f(x)=x^3-ax^2+3x$在$x=1$处取得极值。

（1）求$a$的值；

（2）求$f(x)$在$[-2,2]$上的最大值和最小值。

\item （10分）已知函数$f(x)=\ln x-ax+\dfrac{a}{x}-1$（$a\in\mathbb{R}$）。

（1）当$a=1$时，求$f(x)$的单调区间；

（2）若$f(x)\leq 0$对$x>0$恒成立，求$a$的取值范围。

\item （12分）设函数$f(x)=e^x-ax-1$（$a\in\mathbb{R}$）。

（1）当$a=1$时，证明：当$x\geq 0$时，$f(x)\geq 0$；

（2）若$f(x)$有两个零点$x_1,x_2$（$x_1<x_2$），证明：$x_1+x_2>2$。

\item （12分）已知函数$f(x)=x-\sin x$。

（1）证明：当$x>0$时，$f(x)>0$；

（2）利用（1）的结论，证明：当$x>0$时，$\sin x > x-\dfrac{x^3}{6}$；

（3）进一步证明：当$x>0$时，$x-\dfrac{x^3}{6} < \sin x < x$。

\item （13分）已知函数$f(x)=\ln x+\dfrac{1}{x}-1$。

（1）求$f(x)$的单调区间和极值；

（2）证明：$\dfrac{1}{n+1} < \ln\left(1+\dfrac{1}{n}\right) < \dfrac{1}{n}$对一切正整数$n$成立；

（3）利用（2）的结论，证明：$\left(\dfrac{n+1}{n}\right)^n < e < \left(\dfrac{n+1}{n}\right)^{n+1}$。

\item （13分）已知$f(x)=e^x-x^2-ax-1$（$a\in\mathbb{R}$）。

（1）当$a=1$时，求$f(x)$的最小值；

（2）若$f(x)\geq 0$对$x\in\mathbb{R}$恒成立，求$a$的取值范围；

（3）当$a=0$时，证明：$e^x > 1+x+\dfrac{x^2}{2}$对$x>0$成立。
\end{enumerate}

\end{document}
